% §16.3 Correlated Noise Across Queries
%
% Cross-reference resolution status:
%
% DEFINED here:
%   \label{sec:correlated-noise}              -- this section
%   \label{prop:mcdiarmid}                    -- McDiarmid's inequality
%
% ALREADY DEFINED (chapter 4):
%   \Cref{ax:element-indep}                   -- chapter 4 §4.2 Axiom 1
%   \Cref{sec:axiom-1}                        -- chapter 4 §4.2
%
% ALREADY DEFINED (chapter 3):
%   \Cref{thm:composition}                    -- chapter 3 §3.3 composition theorem
%   \Cref{sec:composition-theorem}            -- chapter 3 §3.3
%
% ALREADY DEFINED (chapter 5):
%   \Cref{sec:bloom-from-scratch}             -- chapter 5 §5.2
%   \Cref{eq:bloom-fpr}                       -- chapter 5 §5.2 FPR equation
%   \Cref{ex:ch5-7}                           -- chapter 5 bit-check independence exercise
%
% INTERNAL forward refs (chapter 16):
%   \Cref{sec:adversarial-inputs}             -- §16.4 (Task 5)
%   \Cref{sec:where-framework-helps}          -- §16.5 (Task 6)
%
% §16.3 covers the second failure mode: error events at distinct queries
% are not independent. Recap Axiom 1; give three diagnostic examples
% (repeated identical query, shared hash families, online updates);
% trace the consequences (chapter 3 composition theorem breaks; chapter
% 5 bit-check independence is exactly this case, per ex 7); state the
% right tools (conditional bounds, McDiarmid's inequality, negative
% correlation salvage); close with the two-Bloom-filter joint-FPR
% worked example.

\section{Correlated Noise Across Queries}
\label{sec:correlated-noise}

The second failure mode trades a different axiom. Where
\Cref{sec:varying-rates} watched the per-query rate drift, this
section keeps the rate fixed and watches the independence between
queries collapse. The channel matrix is stable in the sense of
\Cref{sec:axiom-2}; what fails is the assumption that successive
queries are independent draws from that matrix. Two queries about
the same underlying structure may share noise sources, and the joint
distribution of their outcomes no longer factors as the product of
marginals.

\subsection*{The phenomenon}

Three short examples make the failure pattern concrete; each will
return below in the consequences discussion.

The first example is degenerate but instructive: a single Bloom
filter and a single query, repeated. The Bloom filter's response to
$x \notin A$ is determined by the bit pattern at $k$ specific
positions. Once that pattern is fixed, the response is fixed; the
second query for the same $x$ returns the same answer with
probability one. The conditional probability of a $\top$ on the
second query given a $\top$ on the first equals $1$, not the
marginal $\fprate$. Independence between the two events would
require a conditional probability of $\fprate$; the gap between
$\fprate$ and $1$ is the entire failure.

The second example is less degenerate and far more common. Two
Bloom filters $B_A$ and $B_B$ share a hash family
$h_1, \ldots, h_k$ but represent different sets $A$ and $B$. A
query $x \notin A \cup B$ produces two answers: one from each
filter. Each marginal answer follows the per-filter false-positive
distribution. The joint distribution does not factor: if a hash
function $h_j$ maps $x$ to a position $j^*$ that is set in $B_A$,
the same hash value is checked against $B_B$, and the events that
the bit is set in $B_A$ versus in $B_B$ are not independent of the
hash output for $x$. The correlation runs through the shared
hashes.

The third example introduces time. An online classifier updates
its model after each query response. The classification at query
$t$ depends on the responses to queries $1, \ldots, t - 1$ through
the updated model parameters. Sequential queries therefore have a
dependence structure carried by the model's state at each step.
This is conceptually similar to a Markov chain on the model state;
the per-query rate $\fprate$ may be stable on average, but the
conditional distribution of the $t$-th error given the history is
nontrivial.

\subsection*{Which axiom is violated}

\Cref{sec:axiom-1} stated Axiom 1 as the mutual conditional
independence of error events at distinct inputs, given the input
identities. Its full statement is that, conditional on the true
classifications, the channel's per-input errors are jointly
independent. In symbols: for inputs $x_1, \ldots, x_n$ with true
labels $y_1, \ldots, y_n$ and observed labels
$\hat{y}_1, \ldots, \hat{y}_n$,
\[
  \Pr[\hat{y}_1 = a_1, \ldots, \hat{y}_n = a_n \,|\, y_1, \ldots, y_n]
  \;=\;
  \prod_{i = 1}^n \Pr[\hat{y}_i = a_i \,|\, y_i].
\]
This factorization is the basis of every multi-query bound in the
book. When it fails, the joint becomes an inequality (in either
direction, depending on the sign of the correlation), and the bounds
derived from the factored form are no longer guarantees.

The three examples each fail this in a specific way. The
repeated-identical-query case fails completely: the conditional
probability of the second response given the first is degenerate
($0$ or $1$). The shared-hash-family case fails positively for
queries near collision positions: knowing one filter said $\top$
increases the posterior probability that the other did too. The
online-learning case fails through a temporal Markov dependence:
each error event is conditionally independent of distant past
events given the current model state, but not independent of the
recent past.

A clarification matters here. \Cref{ax:element-indep} is the
conditional-independence form; it can fail while
\Cref{ax:block-rates-indep} (Axiom 2's stability claim) holds. The
two Bloom filters of the second example have stable per-filter
rates; the failure is purely in the joint, not in either marginal.
Time-varying rates and correlated noise are independent failure
modes; either can occur without the other.

\subsection*{Consequences for the iid tools}

The composition theorem of \Cref{sec:composition-theorem}
($\eta_{\text{total}} = 1 - \prod_i (1 - \eta_i)$) used Axiom 1 at
the moment it wrote the product. The product of marginal
no-error probabilities equals the joint no-error probability only
under independence; under positive correlation, the joint is
larger (better than the product predicts), and under negative
correlation, the joint is smaller (worse). The composition
theorem becomes an approximation whose direction of error depends
on the sign of the correlation. The product remains an upper bound
on the joint error rate under positive correlation but is no
longer the joint error rate itself.

The Bloom filter derivation of \Cref{sec:bloom-from-scratch}
contains the same approximation explicitly. The standard
$\fprate \approx (1 - e^{-kn/m})^k$ formula
(\Cref{eq:bloom-fpr}) treats the $k$ bit-check events as
independent: a query $x \notin A$ produces $k$ bit reads
$h_1(x), \ldots, h_k(x)$, and the event ``all $k$ bits are set'' is
written as the product of the per-bit probabilities.
\Cref{ex:ch5-7} of chapter 5 acknowledged that this factorization
is an approximation. The bit positions are negatively correlated
(knowing one bit is set reduces, slightly, the expected number of
items that also set another specific bit), and the true joint
probability is therefore smaller than the product, by an amount
that vanishes as $m \to \infty$ with $n/m$ fixed but is detectable
at small $m$. This is the same Axiom 1 failure, scoped to within a
single query rather than across queries.

The PPV machinery of chapter 7 \S7.4 is less directly affected
because the PPV formula is a per-query quantity. The issue arises
when one combines PPV estimates across queries; the variance of
the empirical PPV depends on the joint distribution of the
contingency-table counts, which under Axiom 1 failure no longer
satisfies the iid sampling assumptions used in
\Cref{prop:ppv-variance}.

\subsection*{The right tools}

Three analytical replacements are standard in the dependent-data
literature, and each has its scope.

\emph{Conditional bounds via the law of total probability.} When
the dependence has a known structure, one can condition on the
relevant variable and bound each conditional separately. For the
two-Bloom-filter case with shared hashes, conditioning on the
hash outputs $h_1(x), \ldots, h_k(x)$ makes the two filter
responses conditionally independent (given the hash positions and
the bit patterns of $B_A$ and $B_B$, the responses are
deterministic functions of the bit patterns, hence trivially
``independent'' given the conditioning). The marginal joint
distribution is then obtained by averaging over the hash output
distribution. This is the most concrete approach when the source
of dependence is structural and identifiable.

\emph{McDiarmid's inequality} for bounded-difference functions of
possibly dependent variables. Where Azuma-Hoeffding required a
martingale structure, McDiarmid bounds a function
$f(X_1, \ldots, X_n)$ that satisfies the bounded-differences
property: changing any one $X_i$ changes the value of $f$ by at
most $c_i$. The $X_i$ themselves need not be independent in the
strict iid sense; the inequality is generally stated for
independent $X_i$ but extends, with the same form, to
martingale-difference sequences of the type Azuma covers.

\begin{proposition}[McDiarmid's bounded-differences inequality]
\label{prop:mcdiarmid}
Let $X_1, \ldots, X_n$ be independent random variables, each
taking values in some measurable space, and let $f$ be a
real-valued function of $(X_1, \ldots, X_n)$ satisfying the
bounded-differences property: for each $i$, there is a constant
$c_i$ such that
\[
  \sup_{x_1, \ldots, x_n, x_i'}
  |f(x_1, \ldots, x_i, \ldots, x_n) - f(x_1, \ldots, x_i', \ldots, x_n)|
  \;\le\; c_i.
\]
Then for any $a > 0$,
\[
  \Pr[|f(X_1, \ldots, X_n) - \mathbb{E}[f]| \ge a]
  \;\le\;
  2 \exp\!\left( -\frac{2 a^2}{\sum_{i = 1}^n c_i^2} \right).
\]
\end{proposition}

The inequality is due to McDiarmid (1989). The connection to
correlated-noise applications is that $f$ might be the empirical
false-positive count of a Bloom filter under a fixed query stream;
the $X_i$ are the bit positions in the filter (set or unset by
inserting the indexed item). The bit positions are not strictly
independent (the constraint that a given item sets exactly $k$
bits ties them together at the level of the hash family), but for
many natural applications, the dependence is weak enough that the
bounded-differences form gives useful bounds.

\emph{Negative correlation as a salvage.} When the dependence is
known to be negative, the iid bounds remain valid as upper bounds
on the dependent system's error rate. This is a quiet but
important rescue: a system whose errors are negatively correlated
behaves better than the iid analysis predicts, not worse. The
classic example is sampling without replacement compared to
sampling with replacement; the former has lower variance, and the
Chernoff bound for the with-replacement case continues to bound
the without-replacement case. For Bernoulli applications,
identifying negative correlation lets one keep the iid analysis
and treat it as conservative.

\subsection*{Worked example: two Bloom filters with shared hashes}

Consider two Bloom filters $B_A$ and $B_B$ sharing a hash family
$h_1, \ldots, h_k$, where $k = 3$, sized at $m = 100$ bits each,
storing $n_A = n_B = 20$ items respectively. The per-filter
false-positive rate, from \Cref{eq:bloom-fpr}, is approximately
\[
  \fprate \;=\; (1 - e^{-k n / m})^k
  \;=\; (1 - e^{-0.6})^3
  \;\approx\; 0.097.
\]
For a query $x \notin A \cup B$, the iid prediction for the joint
event ``both filters say $\top$'' would be
$\fprate^2 \approx 0.0094$. The actual joint event, computed by
conditioning on the hash positions, is larger. The reason is
geometric: the three hash positions $h_1(x), h_2(x), h_3(x)$ are
the same in both filters, and a position that is set in both
filters costs only one ``coincidence'' in the joint, not two
independent coincidences.

A direct calculation gives the actual joint probability. For each
hash position $j$, let $p_A = 1 - (1 - 1/m)^{k n_A} \approx 0.45$
denote the probability that position $j$ is set in $B_A$, and
similarly $p_B \approx 0.45$ for $B_B$. The two events ``position
$j$ is set in $B_A$'' and ``position $j$ is set in $B_B$'' are
independent (because $A$ and $B$ are independent of each other
given the hash family), so the probability that position $j$ is
set in both is $p_A p_B \approx 0.20$. The probability that all
three queried positions are set in both filters, treating the
three positions as independent (which they are not exactly, but
the approximation is acceptable here), is
$(p_A p_B)^3 \approx 0.20^3 \approx 0.008$.

Compare this to the iid prediction $\fprate^2 \approx 0.0094$. The
two numbers are close, but the structure differs: the iid
prediction assumed independent draws from a per-filter marginal,
while the actual joint runs through the shared hash positions and
gives a slightly smaller probability than the product (because the
hash positions are drawn jointly, not independently, between the
two filters). The qualitative point is that the joint depends on
the shared structure, and a system designer who treats the two
filters as independent will get a wrong answer.

A more striking variant: take the two filters to be identical
($A = B$, same hashes, same bit pattern). Then the joint
probability of two $\top$ responses to the same query is exactly
the marginal $\fprate$, not $\fprate^2$. The iid prediction
underestimates the joint by a factor of $1 / \fprate \approx 10$.
The correlation between identical-filter responses is total, and
treating the responses as independent gives a wildly optimistic
error budget.

The diagnostic for correlated noise is to ask whether the noise
sources of the queries are demonstrably independent. If the
queries share hash families, share random seeds, query the same
positions in shared memory, or otherwise touch common randomness,
the responses are correlated, and the iid analysis is an
approximation whose direction and magnitude require explicit
calculation. \Cref{sec:where-framework-helps} returns to which
parts of the framework survive these correlations.
